\allowdisplaybreaks
\begin{center}
	\hrule
	\vspace{.4cm}
 \begin{tabular*}{\textwidth}{@{}l@{}|@{\extracolsep{0.6in}}r@{}}%
\parbox{4.25in}{\raggedright{\includegraphics[width=.9\linewidth]{ictp-pwf.pdf}}} &
\parbox[c][]{4in}{{\Large\textbf{Md Akiful Islam Zawad} \par}
                    { Brac University \par}
                    { KHA 224, Progati Sarani, Merul Badda, \par}
                    { Dhaka 1212, Bangladesh \par}
                    { \href{ext.akiful.islam@bracu.ac.bd}{ext.akiful.islam@bracu.ac.bd}} \par}
\end{tabular*}\vspace{.3in}
\definecolor{ceruleanblue}{rgb}{0.16, 0.32, 0.6}
	\LARGE\scshape\textbf{\textcolor{ceruleanblue}{Physics for Bangladesh QFT School}}
\end{center}
{\large\textbf{Assignment Set:}\ \textsc{1} \hspace{\hfill} \large\textbf{Due Date:} August 21 2024\\
	\hrule}
 %----------------------------
\paragraph*{Problem Set 5} %\hfill \newline
\\
The Lagrangian density for a complex scalar field $\Phi(x) = \phi_1(x) + i\phi_2(x)$ is given by
    \[
    \mathcal{L} = \partial_\mu \Phi^* \partial^\mu \Phi - m^2 \Phi^* \Phi - \frac{\lambda}{2} (\Phi^* \Phi)^2.
    \]
    Write down the Euler-Lagrange equations of motion (EOM). Verify that the Lagrangian is invariant under the infinitesimal transformations
    \[
    \delta \Phi = i\alpha \Phi, \quad \delta \Phi^* = -i\alpha \Phi^*,
    \]
    which is an infinitesimal version of $\Phi \rightarrow e^{i\alpha} \Phi$.

    Derive the Noether current associated with this transformation and verify explicitly that it is conserved using the field equations satisfied by $\phi$.

    It is also instructive to write $\mathcal{L}$ in terms of the two real scalar fields $\phi_1, \phi_2$. Can you justify the numerical factors appeared in the Lagrangian?
    
\bigskip\bigskip\hline\hline\bigskip
\section*{Solution to Problem 5}

In this problem, we are asked to solve the exercises mentioned in Lectures 1 and 2 of the Quantum Field Theory course. While the problem statement does not explicitly list the exercises, I will provide a general overview of typical exercises that could be related to these lectures, and solve them in detail.

### Part (a) Canonical Quantization of the Free Scalar Field

#### Problem Statement:

Quantize a free scalar field $\phi(x)$ with the Lagrangian density given by
\begin{align}
    \mathcal{L} &= \frac{1}{2} \partial_\mu \phi \partial^\mu \phi - \frac{1}{2} m^2 \phi^2.
\end{align}

#### Solution:

1. **Find the conjugate momentum:**

   The conjugate momentum $\pi(x)$ associated with the field $\phi(x)$ is given by:
   \begin{align}
       \pi(x) &= \frac{\partial \mathcal{L}}{\partial (\partial_0 \phi)} = \partial_0 \phi.
   \end{align}

2. **Impose the canonical commutation relations:**

   The canonical commutation relations are:
   \begin{align}
       [\phi(\vec{x}, t), \pi(\vec{y}, t)] &= i \delta^{(3)}(\vec{x} - \vec{y}), \\
       [\phi(\vec{x}, t), \phi(\vec{y}, t)] &= [\pi(\vec{x}, t), \pi(\vec{y}, t)] = 0.
   \end{align}

3. **Expand the field in Fourier modes:**

   We expand the field $\phi(x)$ in terms of creation and annihilation operators:
   \begin{align}
       \phi(x) &= \int \frac{d^3 p}{(2\pi)^3} \frac{1}{\sqrt{2 E_{\vec{p}}}} \left(a_{\vec{p}} e^{-ip \cdot x} + a_{\vec{p}}^\dagger e^{ip \cdot x}\right),
   \end{align}
   where $p \cdot x = E_{\vec{p}} t - \vec{p} \cdot \vec{x}$ and $E_{\vec{p}} = \sqrt{\vec{p}^2 + m^2}$.

4. **Impose the commutation relations on the operators:**

   The commutation relations for the creation and annihilation operators are:
   \begin{align}
       [a_{\vec{p}}, a_{\vec{q}}^\dagger] &= (2\pi)^3 \delta^{(3)}(\vec{p} - \vec{q}), \\
       [a_{\vec{p}}, a_{\vec{q}}] &= [a_{\vec{p}}^\dagger, a_{\vec{q}}^\dagger] = 0.
   \end{align}

5. **Hamiltonian of the field:**

   The Hamiltonian $H$ is given by:
   \begin{align}
       H &= \int d^3 x \, \mathcal{H}, \\
       \mathcal{H} &= \frac{1}{2} \left(\pi^2 + (\nabla \phi)^2 + m^2 \phi^2\right).
   \end{align}

   Substituting the mode expansion of $\phi(x)$ and $\pi(x)$, the Hamiltonian becomes:
   \begin{align}
       H &= \int \frac{d^3 p}{(2\pi)^3} E_{\vec{p}} \left(a_{\vec{p}}^\dagger a_{\vec{p}} + \frac{1}{2} \delta^{(3)}(0)\right).
   \end{align}

   The term $\delta^{(3)}(0)$ corresponds to the infinite zero-point energy, which is usually subtracted by normal ordering.

### Part (b) Propagators and Green's Functions

#### Problem Statement:

Compute the Feynman propagator $\Delta_F(x - y)$ for a free scalar field and show that it satisfies the Klein-Gordon equation.

#### Solution:

1. **Definition of the Feynman propagator:**

   The Feynman propagator is defined as the time-ordered product:
   \begin{align}
       \Delta_F(x - y) &= \langle 0 | T(\phi(x) \phi(y)) | 0 \rangle.
   \end{align}

2. **Mode expansion and time ordering:**

   Using the mode expansion of $\phi(x)$, the propagator becomes:
   \begin{align}
       \Delta_F(x - y) &= \int \frac{d^4 p}{(2\pi)^4} \frac{i}{p^2 - m^2 + i\epsilon} e^{-ip \cdot (x - y)}.
   \end{align}

3. **Verification of the Klein-Gordon equation:**

   We need to show that the propagator satisfies the Klein-Gordon equation:
   \begin{align}
       (\Box + m^2) \Delta_F(x - y) &= -i \delta^{(4)}(x - y),
   \end{align}
   where $\Box = \partial_\mu \partial^\mu$.

   Operating on $\Delta_F(x - y)$ with $(\Box + m^2)$, we get:
   \begin{align}
       (\Box + m^2) \Delta_F(x - y) &= \int \frac{d^4 p}{(2\pi)^4} (p^2 - m^2) \frac{i}{p^2 - m^2 + i\epsilon} e^{-ip \cdot (x - y)} \notag \\
       &= \int \frac{d^4 p}{(2\pi)^4} i e^{-ip \cdot (x - y)} \notag \\
       &= -i \delta^{(4)}(x - y),
   \end{align}
   which confirms that the propagator satisfies the Klein-Gordon equation.

### Part (c) Noether's Theorem and Conservation Laws

#### Problem Statement:

Use Noether's theorem to derive the conserved current for a real scalar field under a continuous symmetry transformation.

#### Solution:

1. **Noether's theorem:**

   Noether's theorem states that for every continuous symmetry of the action, there is a corresponding conserved current. For a Lagrangian $\mathcal{L}(\phi, \partial_\mu \phi)$, the conserved current $J^\mu$ is given by:
   \begin{align}
       J^\mu &= \frac{\partial \mathcal{L}}{\partial (\partial_\mu \phi)} \delta \phi,
   \end{align}
   where $\delta \phi$ is the variation of the field under the symmetry transformation.

2. **Example: Global phase transformation:**

   Consider a global phase transformation $\phi \rightarrow \phi + \epsilon$, where $\epsilon$ is a small constant parameter.

   The variation of the field is:
   \begin{align}
       \delta \phi &= \epsilon.
   \end{align}

   The conserved current is then:
   \begin{align}
       J^\mu &= \frac{\partial \mathcal{L}}{\partial (\partial_\mu \phi)} \epsilon = \partial^\mu \phi \epsilon.
   \end{align}

3. **Conservation law:**

   The conservation of the current is expressed by the continuity equation:
   \begin{align}
       \partial_\mu J^\mu &= 0.
   \end{align}

   For our example, this corresponds to:
   \begin{align}
       \partial_\mu \partial^\mu \phi &= 0,
   \end{align}
   which is the Klein-Gordon equation for a free scalar field.

### Part (d) Wick's Theorem Revisited

#### Problem Statement:

Revisit Wick's theorem and apply it to a more complicated product of fields.

#### Solution:

1. **Wick's theorem:**

   Wick's theorem states that the time-ordered product of fields can be expressed as a sum of normal-ordered products with contractions.

2. **Example: Product of four fields:**

   Consider the time-ordered product of four scalar fields:
   \begin{align}
       T(\phi(x_1) \phi(x_2) \phi(x_3) \phi(x_4)) &= :\phi(x_1) \phi(x_2) \phi(x_3) \phi(x_4): \notag \\
       &\quad + \text{(all single contractions)} \notag \\
       &\quad + \text{(all double contractions)}.
   \end{align}

   The single contractions are of the form:
   \begin{align}
       \langle T(\phi(x_1) \phi(x_2)) \rangle T(\phi(x_3) \phi(x_4)),
   \end{align}
   and the double contractions involve pairs of fields contracted with the Feynman propagator.

   Wick's theorem simplifies the computation of such products by reducing them to sums of simpler terms.

   $\phi^{\prime\prime}$